
%%%%%%%%%%%%%%%%%%%%%%%%%%%%%%%%%%%%%%%%%%%%%%%%%%%%%%%%%%%%%
\subsubsection{具体分析}

% ==== 第1段：问题类型 + 大致方法 ====
问题一要求...，属于<评价/分类/预测/优化/机理分析>类问题。
本问采用...方法进行求解。

% ==== 第2段：完整步骤串联（首先→接着→然后→最后，不换段）====
首先对...进行...，
接着利用...方法...，
然后构建...模型...，
最后通过...得到...。

% ==== 第3段：引出流程图 ====
具体分析流程如图\ref{fig:analysis1_flow}所示。

% ==== 流程图（只改{}中文字，每框≤7字）====
\begin{figure}[H]
  \begin{center}
    \begin{tikzpicture}[x=1cm, y=0.7cm]
      \node[box] (n11) at (0, 0) {步骤1};
      \node[box] (n12) at (5, 0) {步骤2};
      \node[box] (n13) at (10, 0) {步骤3};
      \node[box] (n22) at (5, -3) {步骤5};
      \node[box] (n21) at (0, -3) {步骤4};
      \node[box] (n23) at (10, -3) {步骤6};
      \node[box] (n31) at (0, -6) {步骤7};
      \node[box] (n32) at (5, -6) {步骤8};
      \node[box] (n33) at (10, -6) {步骤9};
      \draw[arrow] (n11) -- (n12);
      \draw[arrow] (n12) -- (n13);
      \draw[arrow] (n23) -- (n22);
      \draw[arrow] (n22) -- (n21);
      \draw[arrow] (n13) -- ++(0,-1.0) -| (n23);
      \draw[arrow] (n21) -- (n31);
      \draw[arrow] (n31) -- (n32);
      \draw[arrow] (n32) -- (n33);
    \end{tikzpicture}
    \caption{问题一分析流程图}
    \label{fig:analysis1_flow}
  \end{center}
\end{figure}

% ============================================================
\subsubsection{模型准备}

% ==== 引文：简述本节做什么 ====
在进行建模之前，需要...。本节将从数据预处理和算法选择两方面进行准备。

% ==== 数据预处理（仅问题1写此段）====
原始数据来源于...，包含...条记录和...个字段。
经检查发现以下问题：...（缺失值/量纲不统一/异常值/编码问题）。

针对上述问题，按以下步骤进行预处理：
第一步，...（如缺失值处理，具体方法+参数）。
第二步，...（如标准化/归一化，写出公式）。
第三步，...（如特征编码/衍生）。

预处理后，数据规模为...行×...列，各指标统计量如下表所示。

% 预处理对比表（可选）
\begin{table}[H]
  \centering
  \caption{数据预处理前后统计量对比}
  \begin{tabularx}{\textwidth}{CCCCC}
    \toprule
    指标 & 预处理前均值 & 预处理前标准差 & 预处理后均值 & 预处理后标准差 \\
    \midrule
    ... & ... & ... & ... & ... \\
    \bottomrule
  \end{tabularx}
  \label{tab:预处理对比}
\end{table}

综上所述，数据预处理完成，为后续建模提供了清洁的输入数据。

% ==== 算法选择 ====
该问题本质上是一个...问题，需要在...条件下...。
常用求解算法包括A算法、B算法、C算法等。本节从维度1、维度2、维度3三方面进行评估。

\begin{table}[H]
  \centering
  \caption{算法能力对比}
  \begin{tabularx}{\textwidth}{CCCCC}
    \toprule
    算法名称 & 维度1 & 维度2 & 维度3 & 稳定性 \\
    \midrule
    算法A（缩写） & 优 & 良 & 中 & 良 \\
    算法B（缩写） & 良 & 优 & 良 & 优 \\
    算法C（缩写） & 中 & 中 & 优 & 中 \\
    \bottomrule
  \end{tabularx}
  \label{tab:算法对比1}
\end{table}

通过对比可以看出，算法A在维度1方面表现最优，...。
因此本文选择算法A作为核心求解方法。

综上所述，本节完成了数据预处理和算法选型，接下来将建立数学模型并进行求解。
